\section{Saddle faces and causal realization}
\label{sec:saddle-realization}
The statistical saddle and its implementation are different objects.
The former is convex; the latter need not be. The purpose of this section
is to retain that distinction while extracting memory obstructions from
an unrestricted saddle. In particular, a zero Bayes coefficient must not
be discarded before the contact conditions at the worst parameters have
been imposed.

\subsection{A normal form over a finite causal interface}
Fix a finite horizon and finite action and report alphabets. An action
$a_t$ is chosen before report $y_t$. A terminal output is included as one
last action, followed by a dummy report. Illegal actions can be removed
from each protocol type; equivalently, use a fixed padding of the finite
protocol tree and impose its zero constraints. Define a \emph{causal
response array} $C(\mathbf a\mid\mathbf y)$ by nonnegativity,
normalization in $\mathbf a$ for every $\mathbf y$, and the requirement
that its marginal on $a_1,\ldots,a_t$ depends only on
$y_1,\ldots,y_{t-1}$. The resulting set $\mathcal C$ is a finite
polytope. Conditioning successive marginals realizes every such array
by a full-history randomized controller; arbitrary versions on zero
marginals do not change the array.

For a fixed parameter $\theta$, the environment has a history-dependent
likelihood $L_\theta(\mathbf a,\mathbf y)$, obtained by marginalizing
its hidden states. The bounded reward has the linear form
\begin{equation}\label{eq:strategic-payoff}
 g_\theta(C)=\sum_{\mathbf a,\mathbf y}
 L_\theta(\mathbf a,\mathbf y)r_\theta(\mathbf a,\mathbf y)
 C(\mathbf a\mid\mathbf y).
\end{equation}
Assume its coefficient vector is continuous on a compact metric
parameter space. The parameter is fixed once, not selected after a
report. Protocol-specific equalities, including a frozen validation
choice, are represented by the corresponding finite protocol tree.
A discarded protocol choice is not supplied again as a free input.

Here is an explicit positive realization of a response array. At cut
$t$ use $k_t$ states, an initial distribution $\alpha$, and nonnegative
matrices $M_t^{a,y}$ of size $k_{t-1}$ by $k_t$ such that
\begin{equation}\label{eq:normalized-instruments}
 M_t^{a,y}\mathbf1=p_t(a),\qquad
 \sum_a p_t(a)=\mathbf1,\qquad
 p_t(a)\text{ is independent of }y.
\end{equation}
Then
\begin{equation}\label{eq:positive-realization}
 C(\mathbf a\mid\mathbf y)
   =\alpha M_1^{a_1,y_1}\cdots M_T^{a_T,y_T}\mathbf1.
\end{equation}
The independence in \eqref{eq:normalized-instruments} is essential:
otherwise an action could use its not-yet-observed report. If an action
or an input must survive an intermediate cut, split this formula at that
cut and include its buffer in the state alphabet. Likewise, impose any
declared shared-row constraints before optimizing. An externally fixed
clock permits time-dependent matrices; an autonomous implementation has
to store its phase or use the prescribed shared rows.

\begin{proposition}[Outer normal form and compatible continuations]
\label{prop:outer-normal-form}
For the finite interface and declared cuts, the response arrays of all
controllers with profile at most $\mathbf K$ are exactly the arrays
\eqref{eq:positive-realization} with $k_t\le K_t$, the normalization
\eqref{eq:normalized-instruments}, and the declared interface constraints.
This set, denoted $\mathcal C_{\mathbf K}$, is compact. The assertion is
over all wirings of the allowed register, not one chosen transition graph.

Equivalently, there are at most $K_t$ normalized continuation arrays
$B_{t,s}$ at each cut, with terminal continuation $1$, and common
nonnegative coefficients satisfying
\begin{equation}\label{eq:continuation-shift}
 B_{t-1,s}(a\mathbf a\mid y\mathbf y)
   =\sum_{s'}M_t^{a,y}(s,s')B_{t,s'}(\mathbf a\mid\mathbf y),
 \qquad C=\sum_s\alpha(s)B_{0,s}.
\end{equation}
Separate positive factorizations at the cuts are not sufficient without
these common shift identities.
\end{proposition}
\begin{proof}
In a controller, let $p_t(a\mid s)$ be its action row and
$q_t(s'\mid s,a,y)$ its update row. Their product defines $M_t^{a,y}$
and gives \eqref{eq:normalized-instruments}. Multiplication over a
complete execution gives \eqref{eq:positive-realization}. Every finite
wiring embeds into the fully connected register by assigning zero to
unused transitions; smaller alphabets are padded by unused states.
Thus no further maximization over graph presentations is necessary.

Conversely choose action $a$ with probability $p_t(a\mid s)$, and,
after observing $y$, use the row
$M_t^{a,y}(s,s')/p_t(a\mid s)$ when the denominator is positive.
At a zero denominator the entire corresponding nonnegative row vanishes;
choose an arbitrary legal update there. This gives exactly the required
runtime arguments. It does not expose a posterior, the parameter, or an
offline occupation array to the machine. Apply the same construction at
each declared buffer cut. A product of finitely many stochastic row
simplices is compact; closed legality and sharing constraints preserve
compactness. Its continuous image is $\mathcal C_{\mathbf K}$.

The future execution starting at a state $s$ defines $B_{t,s}$ and
satisfies \eqref{eq:continuation-shift}. Conversely, induction from the
terminal continuation and these identities gives the product formula.
Thus the continuation description has exactly the same realizations.
In particular, independently selected factorizations at different cuts
cannot be spliced unless the shift equations hold.
\end{proof}

This is a finite controlled positive-realization statement. The role of
positivity and invariant generating cones in stochastic realization is
classical \cite{Heller,Vidyasagar,GG}. Ordinary linear rank need not be a
state count. Here the additional statistical question is which positive
realizations meet an optimal saddle face. Proposition~\ref{prop:outer-normal-form}
provides a normal form, not an efficient factorization algorithm. It
also does not compactify an unbounded number of observable operations.

\subsection{The realization theorem}
Let
\[
 U=\max_{C\in\mathcal C}\min_{\theta}g_\theta(C),\qquad
 V_{\mathbf K}=\max_{C\in\mathcal C_{\mathbf K}}
                                      \min_\theta g_\theta(C).
\]
By compact convex minimax there is a least-favorable prior $\mu_*$ with
$U=\max_{C\in\mathcal C}\int g_\theta(C)d\mu_*$. Put
\begin{align}
 \ell_*(C)&=\int g_\theta(C)d\mu_*,\notag\\
 F_*&=\{C\in\mathcal C:\ell_*(C)=U\},\notag\\
 \mathcal S_*&=F_*\cap
             \bigcap_{\theta}\{C:g_\theta(C)\ge U\}.\label{eq:saddle-face}
\end{align}
The set $F_*$ is an exposed face. The second intersection is important:
a Bayes best response to a least-favorable prior need not be minimax.

Here $\mathcal C$ is the full-history convex reference for the finite
protocol. A nonconvex restriction, such as reuse of one stochastic table
at several visits, is imposed on $\mathcal C_{\mathbf K}$, not inserted
into $\mathcal C$ before invoking convex minimax. In that case $U$ is a
reference upper value unless exactness for the unrestricted physical
class is separately known. In the reset applications below that
exactness is supplied by Theorems~\ref{thm:dual} and \ref{thm:outer}.


\begin{proposition}[Bayes loss and worst-case loss]
\label{prop:loss-decomposition}
Under the preceding assumptions,
\begin{equation}\label{eq:resource-loss}
 U-V_{\mathbf K}=\min_{C\in\mathcal C_{\mathbf K}}
 \left\{U-\ell_*(C)+
                  \ell_*(C)-\min_\theta g_\theta(C)\right\}.
\end{equation}
Both summands are nonnegative and the minimum is attained.
\end{proposition}
\begin{proof}
The prior average is at most $U$ and at least the minimum over fixed
parameters. The sum in braces equals $U-\min_\theta g_\theta(C)$.
Minimizing on the compact realization set proves the assertion.
\end{proof}

\begin{theorem}[Saddle-face contact and compatible realization]
\label{thm:saddle-realization}
Under the preceding assumptions, $V_{\mathbf K}=U$ if and
only if some array in $\mathcal S_*$ has compatible positive
continuation generators of profile at most $\mathbf K$.
If no such generators exist, $V_{\mathbf K}<U$ strictly.

For $C\in\mathcal S_*$ one has $g_\theta(C)=U$ on
$\supp\mu_*$. At an interior supporting parameter at which the payoff
is differentiable, $\nabla_\theta g_\theta(C)=0$. These equalities and
stationarity conditions hold also in directions whose Bayes coefficient
vanishes.

At any reset continuation cut, let $R_h$ denote the conditional channel
of future outputs under each externally specified future input word.
It factors as
\begin{equation}\label{eq:cut-channel-factor}
                  R_h=\sum_{s=1}^{K_t}E(s\mid h)B_s,
 \qquad E(\cdot\mid h)\in\Delta_{K_t},
\end{equation}
with one common family of continuation channels $B_s$.
Any obstruction to this positive factorization for every member of
$\mathcal S_*$ is therefore a memory obstruction for every admissible
architecture, regardless of its training wiring.
\end{theorem}
\begin{proof}
By Proposition~\ref{prop:loss-decomposition}, the two nonnegative
losses vanish exactly when $C$ belongs to $\mathcal S_*$.
Proposition~\ref{prop:outer-normal-form} then proves the realization
criterion. Compactness makes nonintersection
a strictly positive value loss, not merely failure of one optimizer.

For an optimal array, $g_\theta(C)-U$ is continuous, nonnegative, and
has zero $\mu_*$ integral. It is zero on the support of that measure.
Every interior zero is a local minimum, so its gradient vanishes.
No strict-sign hypothesis on individual coordinates is needed.
Finally condition on the actual retained state at a reset cut. Given
that state, the future program receives the same specified inputs and
has the same output law, irrespective of the discarded past. Averaging
gives \eqref{eq:cut-channel-factor}. All persistent random variables
are part of that state. The assertion can equivalently be checked on
controller input tapes, using the generators in
\eqref{eq:continuation-shift}; it does not assume that a hidden
physical state resets at a cut where the interface has not reset it.
\end{proof}

The identity \eqref{eq:resource-loss} is an exact loss decomposition, not
a new interchange of a minimum and maximum over nonconvex machines.
Its use is that an unrestricted certificate identifies a finite face,
while contact at its supporting parameters can determine the remaining
coordinates. The resulting channels are then tested for causal
realizability. This is stronger than choosing an arbitrary Bayes tie rule
and counting the states of that one rule.

\begin{corollary}[Explicit coefficient form for frozen validation]
\label{cor:coefficient-face}
In the reset experiment of \eqref{eq:selector}, write
\[
 C^*_{z\omega}=\int P_\theta^{\otimes N}(z)
                  (Q(\omega)-P_\theta(\omega))\,d\mu_*(\theta).
\]
The first summand of \eqref{eq:resource-loss} is exactly
\begin{equation}\label{eq:weighted-face-loss}
 \sum_{C^*_{z\omega}>0}C^*_{z\omega}(1-h_z(\omega))
 +\sum_{C^*_{z\omega}<0}|C^*_{z\omega}|h_z(\omega).
\end{equation}
Thus optimality fixes every strict-sign coordinate, while the support
contact and stationarity equations must still be imposed at the zeros.
For $K$ labels at the post-training decision cut, an arbitrary upstream
encoder gives the compact relaxation
\begin{equation}\label{eq:decision-factor}
 h_z(\omega)=\sum_{s=1}^K E(s\mid z)D_s(\omega),\qquad
 E\in(\Delta_K)^{\Omega^N},\quad D_s\in[0,1]^\Omega.
\end{equation}
If this relaxation misses the saddle set, its strictly positive gap
bounds every smaller architectural class.
\end{corollary}
\begin{proof}
Maximizing each coefficient separately gives the sum of its positive
part. Subtraction gives \eqref{eq:weighted-face-loss}. Conditioning on
the decision state gives \eqref{eq:decision-factor}; the post-cut
randomization is independent of the discarded word. Conversely an
arbitrary encoder and a response list realize this factorization when
other memory is unrestricted. The parameter spaces in
\eqref{eq:decision-factor} are compact, so the final assertion follows
from Theorem~\ref{thm:saddle-realization}.
\end{proof}

\subsection{Faces as memory witnesses}
For a finite input set $\mathcal Y$ and output alphabet $\mathcal D$,
continuation channels form the polytope
$\mathcal P=\prod_{y\in\mathcal Y}\Delta(\mathcal D)$.
Its vertices are deterministic output functions. The analogous
polytope for scalar test responses is $[0,1]^{\mathcal Y}$.

\begin{lemma}[Vertex and face packing]\label{lem:face-packing}
Suppose channels to be factored through one common state alphabet include
$s$ distinct vertices $v_1,\ldots,v_s$ of a polytope $\mathcal P$ and
rows $r_1,\ldots,r_j$. Suppose $r_i\in F_i$, the faces $F_i$ are
pairwise disjoint, and
\[
 r_i\notin\conv\{v_a:v_a\in F_i\}.
\]
Every factorization of these rows as stochastic mixtures of generators
in $\mathcal P$ needs at least $s+j$ generators.
\end{lemma}
\begin{proof}
A convex decomposition of a vertex uses that same vertex at every
positive coefficient. Hence all $s$ vertices must occur among the
generators. If a convex combination belongs to a face, all generators
with positive coefficients belong to that face. The stated exclusion
therefore forces an additional generator in each $F_i$. Pairwise
disjointness prevents any such generator from serving two faces.
\end{proof}

The lemma does not assert that a stochastic state must itself be a
history residual. Its generator can be any channel in the ambient
polytope. This is why the argument also excludes randomized encoders
and randomized continuations. It is a positive-factorization argument,
not ordinary matrix rank or a restriction to a deterministic compiler.

\begin{proposition}[Resource equivalence and stable obstructions]
\label{prop:saddle-transport}
A causal, cut-preserving recoding with a causal inverse, both using only
one transducer state and preserving the payoff, preserves every
$V_{\mathbf K}$, $U$, and the set of profiles attaining $U$.
More generally suppose simulations in both directions have defect at
most $\epsilon$ and require at most $R_t$ states at cut $t$. For rewards
in $[-1,1]$,
\[
 V_{\mathfrak E}(\mathbf R\mathbf K)
       \ge V_{\mathfrak F}(\mathbf K)-2\epsilon,\qquad
 V_{\mathfrak F}(\mathbf R\mathbf K)
       \ge V_{\mathfrak E}(\mathbf K)-2\epsilon.
\]
For a common controller class and training--validation kernels whose
values differ by at most $\eta$, a gap $U-V_{\mathbf K}=\delta$
remains at least $\delta-2\eta$ after the perturbation.
\end{proposition}
\begin{proof}
Transport a controller through each simulator and retain the product of
its state and the simulator state. Apply the uniform bounded-reward
variation bound and optimize. With an exact one-state inverse the two
inequalities are equalities at the same profile. In the last assertion
both the unrestricted and the constrained optimum change by at most
$\eta$; subtract their two bounds. The state products must be applied
at the actual buffer cuts, not only at the terminal decision.
\end{proof}

These statements place the saddle face, the future-test quotient, and
resource-aware causal morphisms in one theorem chain. The channel
factorizations are presentation-independent under the stated causal
recodings. The memory obstruction in the next section is derived from
the experimental coefficients themselves, not from a metric dimension
or a covering exponent supplied as an additional hypothesis.
